% version 26.05.2011 BK
%\documentclass[showpacs,preprintnumbers,amsmath,amssymb,floatfix]{revtex4}


%\RequirePackage[2020-02-02]{latexrelease} 
% Fixed "extra \endgroup" error

\documentclass[twocolumn,showpacs,preprintnumbers,amsmath,amssymb,floatfix]{revtex4-1}
\usepackage{graphicx}
\usepackage{epsfig}
\usepackage{amsfonts}
\usepackage{amssymb}
\usepackage{epsf}
%\usepackage{epstopdf}
%\newcommand{\insertplot}[5]{\begin{figure}
% \hfill\hbox to 0.05in{\vbox to #5in{\vfill
% \inputplot{#1}{#4}{#5}}\hfill}
% \hfill\vspace{-.1in}
%% \caption{#2}\label{#3}
% \end{figure}}
%\newcommand{\inputplot}[3]{\special{ps:: gsave #2 -36 mul 0 rmoveto
%   currentpoint translate #2 7.0 div #3 5.0 div scale}
% \special{ps: plotfile #1}\special{ps:: grestore}}

%\newcounter{fig}   \newcommand{\lbfig}[1]{\refstepcounter{fig}
%\label{#1} }

%\newcommand{\vphi}{\varphi}

\textwidth=170true mm \textheight=220true mm
\oddsidemargin=-5true mm\evensidemargin=-5true mm
\topmargin=-15true mm

%\def\oneone{\rlap 1\mkern4mu{\rm l}}
%\def\ep{{\epsilon}} 



\usepackage{comment}
\usepackage{color}
\usepackage[normalem]{ulem}
\usepackage{xcolor}
%\newcommand{\fc}[1]{{\textcolor{red}{{#1}}}}
\newcommand{\jk}[1]{{\textcolor{blue}{{#1}}}}
\newcommand{\jl}[1]{{\textcolor{purple}{{#1}}}}
\newcommand{\fsk}[1]{{\textcolor{orange}{{#1}}}}



\begin{document}

\title{
Are slowly rotating Ellis-Bronnikov wormholes stable?}

%\author{ 
%{\bf Humans}
\author{{Bahareh Azad$^1$}}
\email{bahareh.azad@uni-oldenburg.de}
%\affiliation{Institute of  Physics, University of Oldenburg, D-26111 Oldenburg, Germany}
\author{Jose Luis Bl\'azquez-Salcedo$^2$}
\email{jlblaz01@ucm.es}
%\affiliation{Departamento de F\'isica Te\'orica and IPARCOS, Facultad de Ciencias F\'isicas, Universidad Complutense de Madrid, Spain}
\author{Fech Scen Khoo$^1$}
\email{fech.scen.khoo@uni-oldenburg.de}
%\affilation{Institute of  Physics, University of Oldenburg, D-26111 Oldenburg, Germany}
\author{{Jutta Kunz$^1$}}
\email{jutta.kunz@uni-oldenburg.de} 
%\affilation{Institute of  Physics, University of Oldenburg, D-26111 Oldenburg, Germany}


\affiliation{
$^1$Institut f\"ur  Physik, Universit\"at Oldenburg, Postfach 2503,
D-26111 Oldenburg, Germany
\\
$^2$Departamento de F\'isica Te\'orica and IPARCOS, Facultad de Ciencias F\'isicas, Universidad Complutense de Madrid, Spain
}
\date{\today}
\pacs{}
\begin{abstract}
We investigate the radial perturbations of Ellis-Bronnikov wormholes ($\mathrm{l}=0$) in a slowly rotating background expanded up to second order in rotation.
We find indications that simple wormhole solutions such as Ellis-Bronnikov in General Relativity can be stabilized by rotation, thus favoring a viable traversable wormhole. 
This opens up the intriguing question whether the many other wormhole solutions with or without the support of exotic matter can become linearly mode stable when the wormhole rotates.
\end{abstract}

\maketitle

%\begin{document}

%%%%%%%%%%%%%%%%%%%%%%%%%%%%%%%%%%%%%%%%%%%%%%%%%%%%%%%%%%%%%%%%%%%%%%%%%%%%%%
%%%%%%                     Introduction
%%%%%%%%%%%%%%%%%%%%%%%%%%%%%%%%%%%%%%%%%%%%%%%%%%%%%%%%%%%%%%%%%%%%%%%%%%%%%%
\noindent{\textbf{Introduction.}}
%
Wormholes have received much attention in recent years, in particular, concerning their possible astrophysical signatures such as gravitational lensing effects of wormholes
\cite{Cramer:1994qj,Safonova:2001vz,Perlick:2003vg,Nandi:2006ds,Abe:2010ap,Toki:2011zu,Nakajima:2012pu,Tsukamoto:2012xs,Kuhfittig:2013hva,Bambi:2013nla,Takahashi:2013jqa,Tsukamoto:2016zdu},
wormhole shadows
\cite{Bambi:2013nla,Nedkova:2013msa,Ohgami:2015nra,Shaikh:2018kfv,Gyulchev:2018fmd,Guerrero:2022qkh},
or wormhole accretion disks and quasi-periodic oscillations
\cite{Harko:2008vy,Harko:2009xf,Bambi:2013jda,Zhou:2016koy,Lamy:2018zvj,Deligianni:2021ecz,Deligianni:2021hwt}.
In fact, wormholes may mimic many properties of black holes 
\cite{Damour:2007ap,Bambi:2013nla,Azreg-Ainou:2014dwa,Dzhunushaliev:2016ylj,Cardoso:2016rao,Konoplya:2016hmd,Nandi:2016uzg,Bueno:2017hyj,Blazquez-Salcedo:2018ipc,Azad:2020ajs,Azad:2022qqn}.
Therefore the study of gravitational waves emitted by wormholes is also of great interest.

The observations of gravitational waves by the LIGO/VIRGO collaboration have provided us with invaluable information about the merger of highly compact objects such as black holes and neutron stars (see e.g. \cite{LIGOScientific:2016aoc,LIGOScientific:2017vwq}). 
Here the ringdown phase after the merger events is dominated by the characteristic frequencies of the final objects, encoded in their quasi-normal modes (QNMs).
Whereas the QNMs of black holes and neutron stars are well studied in General Relativity (GR) \cite{Kokkotas:1999bd,Berti:2009kk,Konoplya:2011qq}, the corresponding modes of wormholes have received much less attention. % 

In GR classical wormholes need exotic matter in order to violate the energy conditions \cite{Ellis:1973yv,Ellis:1979bh,Bronnikov:1973fh,Morris:1988tu,Morris:1988cz,Visser:1995cc,Hochberg:1996ee,Hochberg:1998ii,Visser:2003yf,Bronnikov:2005gm,Alcubierre:2017pqm}.
%
While this issue can be circumvented in several ways (for example, by considering wormholes in generalized gravities \cite{Kanti:2011jz,Alcubierre:2017pqm} or by adding fermionic matter \cite{Blazquez-Salcedo:2020czn,Konoplya:2021hsm}),
%
another question concerns the stability of wormholes in GR and generalized gravities. 
For instance, the notorious radial instability of the static spherically symmetric Ellis-Bronnikov (EB) wormholes of GR has been known for long \cite{Shinkai:2002gv,Gonzalez:2008wd,Gonzalez:2008xk,Cremona:2018wkj,Blazquez-Salcedo:2018ipc}.


Here we present intriguing evidence that this radial instability of EB wormholes will disappear, when these wormholes rotate sufficiently fast.
However, unlike their static counterparts, rotating EB wormholes are known only perturbatively for slow rotation \cite{Kashargin:2007mm,Kashargin:2008pk} and numerically for fast rotation \cite{Kleihaus:2014dla,Chew:2016epf}.
Additionally, rotating EB wormholes have been investigated in five dimensions for the case of two equal magnitude angular momenta, where the angular coordinates factorize \cite{Dzhunushaliev:2013jja}.

In four dimensions, a mode analysis of rotating compact objects becomes a challenging task.
Already for Kerr-Newman black holes linear mode stability could only be provided by solving the resulting system of coupled partial differential equations with sophisticated numerical methods \cite{Dias:2015wqa}.
Alternatively, the modes have recently been obtained perturbatively to second order in rotation, where parity mixing complicates their evaluation as opposed to first order \cite{Blazquez-Salcedo:2022eik}.

We here apply this recently developed second order formalism \cite{Blazquez-Salcedo:2022eik} in order to study the change of the unstable radial mode of EB wormholes in the presence of slow rotation.
We demonstrate, that the magnitude of the pertinent unstable eigenvalue decreases with increasing angular momentum of the wormhole, thus diminishing the instability.
In fact, in our perturbative analysis the eigenvalue 
%
{vanishes}
at some critical value of the angular momentum, that slightly depends on the asymmetry parameter of the wormhole, {and we argue that nonperturbatively the instability will disappear slightly earlier}.


%%%%%%%%%%%%%%%%%%%%%%%%%%%%%%%%%%%%%%%%%%%%%%%%%%%%%%%%%%%%%%%%%%
\noindent{\textbf{EB Wormholes in Slow Rotation.}} 
%%%%%%%%%%%%%%%%%%%%%%%%%%%%%%%%%%%%%%%%%%%%%%%%%%%%%%%%%%%%%%%%%%
We consider a phantom scalar field $\Phi$ minimally coupled to GR,
\begin{eqnarray}
			S &=& \frac{1}{16 \pi G}\int d^4x \sqrt{-g} 
		\Big[\mathrm{R} + 2 \partial_\mu \Phi \, \partial^\mu \Phi 
		 \Big] \, 
   \label{eq:ellis} 
\end{eqnarray}
with field equations %
\begin{eqnarray}
    \mathrm{R}_{\mu\nu} = -2 \partial_{\mu}\Phi\partial_{\nu}\Phi \, , \quad 
    \partial_{\mu}\partial^{\mu}\Phi = 0 \, .
\end{eqnarray}
%
%\newpage
%\noindent 
Up to second order in rotation the background metric is given by
\begin{eqnarray}
ds^2 &=& -e^{f}\left[1+\epsilon_r^2 2\left(h_0(r)+h_2(r)P_2(\theta)\right)\right]dt^2 
\nonumber 
\\ 
%
&+& e^{-f}\left[1+\epsilon_r^2 2\left(b_0(r)+b_2(r)P_2(\theta)\right)\right]dr^2 \nonumber \\ 
&+& 
%
{e^{-f}} R^2
\left[1+\epsilon_r^2 2\left(k_0(r)+k_2(r)P_2(\theta)\right)\right]
\nonumber 
\\ 
&& \times 
\left\{ d\theta^2+\sin^2{(\theta)}\left[d\varphi-\epsilon_rw(r)dt\right]^2 \right\} \ ,
\label{metric_1}
\end{eqnarray}
where $R^2=(r^2+r_0^2)$, with slow-rotation parameter $\epsilon_r \ll 1$, and the Legendre polynomial $P_2(\theta) = \left(3\cos^2{(\theta)}-1\right)/2$.
%
%
%
The corresponding second order phantom scalar field has the form
\begin{eqnarray}
\Phi = \phi(r) + \epsilon_r^2 \left(\phi_{20}(r)+\phi_{22}(r)P_2(\theta)\right) \ .
\label{scalar_1}
\end{eqnarray}
%
In the following we choose the gauge $k_0(r)=0$.
%
The throat of the wormhole is located at $r=0$, and it has an area {$\mathrm{A}=4\pi r_0^2 e^{\frac{\pi C}{2r_0}}$.}
%
%

The EB wormhole background solutions are
\begin{eqnarray}
    f(r) = \frac{{C}}{r_0} \left(\tan^{-1} \! \left(\frac{r}{r_0}\right)-\frac{\pi}{2}\right) \ , \ \phi(r)= \frac{Qf}{C} \ ,
\end{eqnarray}
%
%
with asymmetry parameter $C=2M_0$, related to the static background wormhole mass $M_0$, and the static background phantom scalar charge %
%
%
{$Q_{0} = \sqrt{C^2/4 + r_0^2}$.}
%
%
%
%
%
%
%
%
Here we will focus on $C\ge 0$.
At the first order in rotation, we solve to obtain
%
\begin{eqnarray}
    w(r)= \frac{3J}{2C(C^2+r_0^2)}
    \left[
    1 - \left(
    1+2C\frac{C+r}{R^2}
    \right)e^{2f}
    \right] \, . \, \, \, \, \, \, \, \,
\end{eqnarray}
In the system of field equations, the differential equations of the radial functions $h_0,{b_0},\phi_{20}$ are decoupled from those of ${h_2}, k_2, {b_2}, \phi_{22}$.
%
When focusing on radial perturbations, we only need the functions $\phi_{20}$, {$b_0$}, $h_0$.
For the simpler 
%
{symmetric}
background wormholes ($C=0$)
%
they are given by
\begin{eqnarray}
    \phi_{20}(r) &=& 
    \frac{3J^2}{rr_0^3}
    \left( 
    \phi + \frac{\pi}{2} - \frac{3r}{2r_0} - \frac{4}{\pi r_0} - \frac{3r^2+4r_0^2}{2R^2\phi}
    \right)
    \, , \nonumber \\
%
%
    b_0(r) &=& \frac{\Delta M}{r}
    \left(
    1-\phi\frac{r_0}{r}
    \right)
    +
    \frac{6J^2}{rr_0^3}
    \left(
    \frac{r_0}{r}+\phi
    \right)
    -\frac{3J^2}{r^2r_0^2}\phi^2
        \, , \nonumber \\
%
%
    h_0(r) &=&  \frac{3J^2}{r_0^4}
    \left(
    \phi^2 + \frac{\pi^2-8}{2\pi}\phi - \frac{r_0^2}{R^2}
    \right) \, .
\end{eqnarray}

%
Solutions for $C>0$ are more involved, and hence provided 
%
{elsewhere \cite{Azad:2023}}.
The static mass $M_0$ is corrected by
$\Delta M =3J^2(\pi^2-8)/(2\pi r_0^3)$
%
%
%
in the limit of $C=0$ when $r\rightarrow +\infty$.
Therefore, for %
$C=0$, the total mass of the wormhole is no longer zero. %
In fact, it receives a contribution from the rotation.
This echoes the finding in \cite{Kashargin:2008pk},
that there is no massless rotating {EB} wormhole.
%
The static scalar charge $Q_0$ is also corrected by $\Delta Q=-\Delta M$.
%
The mass {and scalar charge corrections for $C>0$ are} given
%
{elsewhere \cite{Azad:2023}}.

\noindent{\textbf{Stability Analysis:
Radial Perturbations.}} 
%
We linearly perturb the metric field $g$ and the phantom field $\Phi$ up to first order in the perturbation parameter $\epsilon_q$, in our slowly rotating ($sr$) background considered up to second order in rotation $\epsilon_r^2$.
The perturbed metric is then given by
\begin{eqnarray}
g_{\mu\nu} &=& g^{(sr)}_{\mu\nu} + \epsilon_q \delta h_{\mu\nu}(t,r,\theta,{\varphi}) \, , 
\end{eqnarray}
%
%
the perturbation of the metric being a combination of polar ($P$) and axial ($A$) perturbations:
$ \delta h_{\mu\nu}=\delta h^{(A)}_{\mu\nu} + \delta h^{(P)}_{\mu\nu}$.
The perturbed phantom field is
\begin{eqnarray}
\Phi &=& {\Phi}^{(sr)} + \epsilon_q \delta{\phi}^{(P)} (t,r,\theta,{\varphi})
%
\, .
\end{eqnarray}
These perturbations can be decomposed in spherical harmonics 
$Y[\mathrm{l},\mathrm{m}](\theta,{\varphi})$ , 
with multipole numbers $\mathrm{l}$ and $\mathrm{m}$  \cite{Blazquez-Salcedo:2022eik}.
When wormholes rotate the perturbations of even (polar) and odd (axial) parities are entangled, i.e.~the multipole numbers $\mathrm{l}$ are coupled.
This results in an infinite tower of equations carrying the index $\mathrm{l}$ that are to be summed, as is well-known for other rotating objects such as black holes and neutron stars.
%
Here, we are interested in the radial-led perturbations, hence $\mathrm{l}=0$, $\mathrm{m}=0$ polar-led perturbations. 
For a given slowly rotating background up to second order in rotation, it is possible to show that these perturbations are just coupled to axial $\mathrm{l}=1$ perturbations ({see} \cite{Blazquez-Salcedo:2022eik}). Hence for our purposes, the metric perturbation can be simplified {to}
\begin{eqnarray}
\delta h_{\mu\nu}= e^{i\omega t}
    \begin{pmatrix}
NY_0 & 0    & 0 & {S}_0 Y_\theta \\
0   & LY_0 & 0 & {S}_1 Y_\theta \\
0 & 0 & R^2 TY_0 & 0 \\
{S}_0 Y_\theta & {S}_1 Y_\theta & 0 & R^2 T\sin^2{\theta}Y_0
\end{pmatrix} \, \, \, \, \, \, \, \, \,
\end{eqnarray}
%
%
%
%
where 
$N, T, L, {S}_0, {S}_1$ are radial functions, 
$Y_0=Y[0,0]
=1/\sqrt{4\pi}$ 
and 
$Y_\theta=\sin{\theta} \, \partial_\theta Y[1,0]
=-\sin^2{\theta}\sqrt{3/4\pi}$. 
The phantom perturbation on the other hand can be taken as
\begin{eqnarray}
 \delta{\phi}^{(P)} = 
 %
 {\phi_1(r)} Y_0 \, .
\end{eqnarray}
%
%
With these perturbations, a consistent set of equations is formed, as the rotations are approximately slow and we {work} only up to second order in the rotation parameter $\epsilon_r$. More technical details of our slow-rotation approximation framework can be found in \cite{Blazquez-Salcedo:2022eik}.

Furthermore, the gauge can be fixed by taking 
%
${S}_0={\phi_1}=0$ 
to simplify the resulting system of equations. It can be shown that the 
{final}
system can be cast into a single second order ordinary differential equation for $T$, with algebraic relations for the rest of the perturbation functions.
%
%
%
%
If we take $T =\alpha Z$, the perturbation equation can be cast into a Schrödinger equation,
\begin{eqnarray}
%
    \partial^2_{r_*}Z
    + (\omega^2 - V(r))Z = 0 \, ,
\end{eqnarray}
where $\omega$ is the quasi-normal mode frequency, $V(r)$ is the potential, and the tortoise coordinate $r_*$ and amplitude $\alpha$ are given by,
\begin{eqnarray}
%
   \partial_r r_*
   = e^{-f}
   \left[
   1 +
%
   \epsilon_r^2 (e^{-f}b_0-h_0)
   \right] \, ,   \, \, \, \, \, \, \, \, \, \, \, \,  \\
%
   \frac{\partial_r\alpha}{\alpha}
%
   = \frac{(\partial_r\phi)^2R^2
   +\epsilon_r^2
   [
   2e^{-f}b_0
%
   -
%
   e^{-2f}R^4(\partial_r w)^2/6
   ]
%
   }{r-\sqrt{R^2(\partial_r\phi)^2-r_0^2}}
%
   \, .  \, \, \, \, \, \, \, \, \, \, \, \, 
\end{eqnarray}
%
%

In this case, the potential is given by
\begin{eqnarray}
    V(r) &=& V_{static}(r) + J^2 \, 
%
    V_2(r) \, .
%
%
    \label{pot_term}
\end{eqnarray}
There is no linear term in the angular momentum $J$.
Its contribution only enters at second order,
and it is a function of the static background solutions $f$, $\phi$, 
the first order function $w$, and the second order functions $h_0$, $b_0$, $\phi_{20}$.
Taking $C$ to zero gives rise to the following potential
\begin{eqnarray}
&& V_{C=0}\,{(r)} =\frac{3r^2r_0^2+2r_0^4}{r^2R^4}+ 
%
\\
&&
J^2 
\Bigg(
\frac{-12\phi^2}{r^4R^2}+
 \frac{12\pi r
(R^4+r_0^4)
-6r_0^3
(\pi^2-8)
R^2}
{\pi r_0^3 r^4 R^4}\phi+
\nonumber
\\
&& \frac{
3r(\pi^2-8)
%
(R^4+r_0^4)
+12\pi r_0
%
(R^2r^2+2r_0^4)
}{\pi r_0^3 r^4 R^4}
\Bigg) \, . \nonumber
\end{eqnarray}
We provide the complete expression of the potential term for $C>0$ %
{elsewhere \cite{Azad:2023}}.
In Figure \ref{Fig_pot}, we {show} %
the correction to the potential (\ref{pot_term}) of the perturbation system due to slow rotation, obtained with 
%
{symmetric}
($C=0$) and 
%
{asymmetric}
($C\ne 0$) background wormhole configurations.
The potential shown has been scaled with %
%
%
$\left(r-C/2\right)^2$.

\noindent{\textbf{Numerical Method.}} 
Here we are interested in unstable modes, {therefore} %
we focus on perturbations for which the eigenvalue $\omega$ is a purely imaginary number, $\omega = i\omega_I$, such that $\omega_I<0$. 
In this case, an asymptotic analysis of the perturbation function $T(r)$ reveals that it has to decay exponentially to zero at both infinities, $T(\pm\infty)=0$. In practice we solve the second order ODE for the function $T(r)$ together with 
%
an
auxiliary equation $E'=0$ where $E=\omega^2$. 
%
Non-trivial regular solutions are obtained by fixing the amplitude of the perturbation function to an arbitrary real constant at an arbitrary point $r=r_c$, i.e. $T(r_c)=constant$. This method has been used successfully before to generate unstable modes of the static EB wormholes \cite{Blazquez-Salcedo:2018ipc}.
%
%
As an illustrative example, we show in Figure \ref{Fig_pot} a typical Gaussian profile of the function $T(r)$.
%
%
\begin{figure}[h!]
		\centering
    		\includegraphics[width=0.32\textwidth,angle=-90]{figs/plot_pot_c0_v2.eps}
		\caption{Potential correction $\tilde{V}_2=V_2(r-C/2)^2$ {vs radial coordinate $r$} for several values of $C$ for $r_0=1$, {and} %
  a typical profile of the perturbation function $T(r)$ {(in grey)}.} % \vspace{-0.25cm}
		\label{Fig_pot}
	\end{figure}	
%
%
%

%
In order to generate the values of $\omega_I$, we start from the static limit ($J=0$), fixing $r_0=1$ and an arbitrary value of $C$. We reproduce the results from \cite{Blazquez-Salcedo:2018ipc}, and by slowly increasing the value of the angular momentum $J$, we get a relation
\begin{eqnarray}
    \omega_I = \omega_I^{(0)} + J^2 \Delta\omega_I^{(2)} \, ,
\end{eqnarray}
where $\omega_I^{(0)}$ is the unstable mode of the static EB wormhole, and $\Delta\omega_I^{(2)}$ is the slow rotation correction to the mode.




\noindent{\textbf{%
Radial Instability.}} %
We exhibit in Figure \ref{Fig_J_omegaI_inset} the dependence of the unstable radial mode on the scaled and dimensionless angular momentum $J/A$ of the EB wormholes, where $A$ is the area of the throat (obtained in second order).
We recall that rotating EB wormholes approach the extremal Kerr black hole as their limiting configuration for large angular momenta \cite{Kleihaus:2014dla,Chew:2016epf}.
%
%
Thus the limiting case {here} is given {by} $J/A\to 1/(8\pi) \approx 0.0397$. Since {fast} rotating EB wormholes cannot exceed this limit, 
%
%
%
the slow rotation approximation {can} only make sense %
%
%
{below} this limit.
%
%



In Figure \ref{Fig_J_omegaI_inset} we show that, as we increase the value of the angular momentum, the value of $|\omega_I|$ decreases to zero. This is because, for every value of $C$ we have studied, the second order correction to the mode is positive, $\Delta\omega_I^{(2)}>0$. 
%
As a result, it is always possible to obtain a critical value of the angular momentum, $J_c$, for which the %
{eigenvalue reaches zero} %
and hence {this unstable mode}
%
%
disappears. 
{
%
}
Most importantly, {this happens}
%
at a {value of $J$} much smaller %
{than the} limiting value.
%
{In fact,} as seen in the inset of Figure \ref{Fig_J_omegaI_inset}, the 
critical values of the scaled angular momentum {$J_c/A$ decrease monotonically with C.}
{Moreover,} for each of the values of $C$ examined, 
%
{$J_c/A$} is around $50\%$ of the limiting value, which is well within the regime for which the second order slow rotation approximation has been shown to work {well} (for example, the quadratic approximation reproduces with excellent precision the spectrum of QNMs of Kerr-Newman black holes \cite{Blazquez-Salcedo:2022eik}).
%
%
%
%
%
%
%
%
\begin{figure}[h!]
		\centering
    		\includegraphics[width=0.32\textwidth,angle=-90]{figs/plot_J_omegaI_b1_new.eps}
		\caption{{Main} unstable mode {$\omega_I \sqrt{A}$ vs angular momentum $J/A$ (scaled with throat area $A$)} for several values of $C$. Inset: Critical {values $J_c/A$ vs $C$}.}  % \vspace{-0.25cm}
		\label{Fig_J_omegaI_inset}
	\end{figure}	


  \begin{figure}[h!]
		\centering
    		\includegraphics[width=0.32\textwidth,angle=-90]{figs/plot_J_omega2_new4.eps}
		\caption{Main unstable mode {(solid)} %
  and second unstable mode {(dotted)
  vs angular momentum $J/A$ (scaled with throat area $A$)} 
  for $C=0$, $0.05$, $0.1$ and $0.2$.} % \vspace{-0.25cm}
		\label{Fig_J_w2_both}
	\end{figure}	

However, the critical {value} {of the} angular momentum
for which the wormholes {become stable} %
is {presumably still} %
smaller, as the following %
{observations} suggest.
%
{When} the wormholes are set into rotation, a zero mode of the static wormholes changes into a second unstable radial mode of the slowly rotating EB wormholes.
Following a quadratic dependence with a small negative coefficient $\Delta\omega_I^{(2)}$, this mode remains at first very small in magnitude.
Subsequently, its magnitude increases and it crosses the corresponding main unstable mode {not too far from} %
its critical {value $J_c/A$}. %
We exhibit both unstable modes in Figure \ref{Fig_J_w2_both}. 
%
Note that for the values of $C$ displayed the crossing occurs around $J/A\approx0.013$, slightly above $30\%$ of the limiting {value}. %
%

We do not continue both modes beyond their crossing, however, since {this crossing} is indicating a %
{deficiency} of the applied quadratic approximation.
A non-perturbative study will allow for a far stronger dependence on the angular momentum than the current quadratic one.
We therefore conjecture, that when nearing the critical value, both unstable modes will change much more rapidly, approach each other, merge and disappear.
They will thus reflect the behavior observed in higher dimensions in the special case of wormholes with equal angular momenta \cite{Dzhunushaliev:2013jja}.


\noindent{\textbf{Conclusions.}} 
While EB wormholes represent rather simple wormhole solutions they may be considered prototypes whose properties can teach a lot about wormholes in general.
Therefore their response to rotation is highly remarkable.
Indeed, the potential of rotation to stabilize previously unstable wormholes was suggested before (see e.g. \cite{Matos:2005uh}), but we have given first evidence here, that it can indeed happen for wormholes in four space-time dimensions, where rotation implies significant loss of symmetry (associated with the technical challenge).

As we have shown above, the notorious unstable mode of static EB wormholes becomes more stable, as slow rotation is turned on.
In our second order and thus quadratic approximation the mode then becomes stable at a critical value of the angular momentum, which depends on the asymmetry parameter $C$.
For the EB wormholes we have considered, this critical value is around $50\%$ of the limiting angular momentum.
%
Moreover, we have observed that for rotating wormholes a second mode appears, {and} similar to the higher dimensional case \cite{Dzhunushaliev:2013jja}, it tends to merge with the fundamental unstable mode. 
Hence we have conjectured the behavior of the radial instability for the case of a future non-perturbative study: the radial instability will disappear at even smaller values of the angular momentum, around $30\%$ of the limiting {value}. 
%
%
Thus for {sufficiently} rapid rotation the 
%
wormholes will be stable.



%%%%%%%%%%%%%%%%%%%%%%%%%%%%%%%%%%%%%%%%%%%%%%%%%%%%%%%%%%%%%%%%%%%%%%%%%%%%%%%%%%
\noindent{\textbf{Acknowledgements.}}
We gratefully acknowledge support by DAAD, the DFG Research Training Group 1620 \textit{Models of Gravity}, 
DFG project Ku612/18-1, 
FCT project PTDC/FIS-AST/3041/2020, COST Actions CA15117 and CA16104, and MICINN project PID2021-125617NB-I00 ``QuasiMode".
JLBS gratefully acknowledges support from Santander-UCM project PR44/21‐29910.

%%%%%%%%%%%%%%%%%%%%%%%%%%%%%%%%%%%%%%%%%%%%%%%%%%%%%%%%%%%%%%%%%%%%%%%%%%%%%%%%%%






\begin{thebibliography}{99}

%\cite{Cramer:1994qj}
\bibitem{Cramer:1994qj}
  J.~G.~Cramer, R.~L.~Forward, M.~S.~Morris, M.~Visser, G.~Benford and G.~A.~Landis,
  %``Natural wormholes as gravitational lenses,''
  Phys.\ Rev.\ D {\bf 51}, 3117 (1995)
  %doi:10.1103/PhysRevD.51.3117.
%  [astro-ph/9409051].

%\cite{Safonova:2001vz}
\bibitem{Safonova:2001vz}
  M.~Safonova, D.~F.~Torres and G.~E.~Romero,
  %``Microlensing by natural wormholes: Theory and simulations,''
  Phys.\ Rev.\ D {\bf 65}, 023001 (2002)
  %doi:10.1103/PhysRevD.65.023001.
%  [gr-qc/0105070].

%\cite{Perlick:2003vg}
\bibitem{Perlick:2003vg}
  V.~Perlick,
  %``On the Exact gravitational lens equation in spherically symmetric and static space-times,''
  Phys.\ Rev.\ D {\bf 69}, 064017 (2004)
  %doi:10.1103/PhysRevD.69.064017.
 % [gr-qc/0307072].

%\cite{Nandi:2006ds}
\bibitem{Nandi:2006ds}
  K.~K.~Nandi, Y.~Z.~Zhang and A.~V.~Zakharov,
  %``Gravitational lensing by wormholes,''
  Phys.\ Rev.\ D {\bf 74}, 024020 (2006)
  %doi:10.1103/PhysRevD.74.024020.
 % [gr-qc/0602062].

%\cite{Abe:2010ap}
\bibitem{Abe:2010ap}
  F.~Abe,
  %``Gravitational Microlensing by the Ellis Wormhole,''
  Astrophys.\ J.\  {\bf 725}, 787 (2010)
 % doi:10.1088/0004-637X/725/1/787.
  %[arXiv:1009.6084 [astro-ph.CO]].

%\cite{Toki:2011zu}
\bibitem{Toki:2011zu}
  Y.~Toki, T.~Kitamura, H.~Asada and F.~Abe,
  %``Astrometric Image Centroid Displacements due to Gravitational Microlensing by the Ellis Wormhole,''
  Astrophys.\ J.\  {\bf 740}, 121 (2011)
  %doi:10.1088/0004-637X/740/2/121.
 % [arXiv:1107.5374 [astro-ph.CO]].

%\cite{Nakajima:2012pu}
\bibitem{Nakajima:2012pu}
  K.~Nakajima and H.~Asada,
  %``Deflection angle of light in an Ellis wormhole geometry,''
  Phys.\ Rev.\ D {\bf 85}, 107501 (2012)
  %doi:10.1103/PhysRevD.85.107501.
 % [arXiv:1204.3710 [gr-qc]].

%\cite{Tsukamoto:2012xs}
\bibitem{Tsukamoto:2012xs}
  N.~Tsukamoto, T.~Harada and K.~Yajima,
  %``Can we distinguish between black holes and wormholes by their Einstein ring systems?,''
  Phys.\ Rev.\ D {\bf 86}, 104062 (2012)
  %doi:10.1103/PhysRevD.86.104062.
 % [arXiv:1207.0047 [gr-qc]].

%\cite{Kuhfittig:2013hva}
\bibitem{Kuhfittig:2013hva}
  P.~K.~F.~Kuhfittig,
  %``Gravitational lensing of wormholes in the galactic halo region,''
  Eur. Phys.\ J. \ C {\bf 74},  2818 (2014)
  %doi:10.1140/epjc/s10052-014-2818-2.
 % [arXiv:1311.2274 [gr-qc]].
 
  %\cite{Bambi:2013nla}
\bibitem{Bambi:2013nla}
  C.~Bambi,
  %``Can the supermassive objects at the centers of galaxies be traversable wormholes? The first test of strong gravity for mm/sub-mm very long baseline interferometry facilities,''
  Phys.\ Rev.\ D {\bf 87}, 107501 (2013)
  %doi:10.1103/PhysRevD.87.107501.
  %[arXiv:1304.5691 [gr-qc]]

%\cite{Takahashi:2013jqa}
\bibitem{Takahashi:2013jqa}
  R.~Takahashi and H.~Asada,
  %``Observational Upper Bound on the Cosmic Abundances of Negative-mass Compact Objects and Ellis Wormholes from the Sloan Digital Sky Survey Quasar Lens Search,''
  Astrophys.\ J.\  {\bf 768}, L16 (2013)
  %doi:10.1088/2041-8205/768/1/L16.
%  [arXiv:1303.1301 [astro-ph.CO]].

%\cite{Tsukamoto:2016zdu}
\bibitem{Tsukamoto:2016zdu}
  N.~Tsukamoto and T.~Harada,
  %``Light curves of light rays passing through a wormhole,''
  Phys.\ Rev.\ D {\bf 95},   024030 (2017)
  %doi:10.1103/PhysRevD.95.024030.
 % [arXiv:1607.01120 [gr-qc]].

%\cite{Nedkova:2013msa}
\bibitem{Nedkova:2013msa}
  P.~G.~Nedkova, V.~K.~Tinchev and S.~S.~Yazadjiev,
  %``Shadow of a rotating traversable wormhole,''
  Phys.\ Rev.\ D {\bf 88},   124019 (2013)
  %doi:10.1103/PhysRevD.88.124019.
%  [arXiv:1307.7647 [gr-qc]].

%\cite{Ohgami:2015nra}
\bibitem{Ohgami:2015nra}
  T.~Ohgami and N.~Sakai,
  %``Wormhole shadows,''
  Phys.\ Rev.\ D {\bf 91}, 124020 (2015)
  %doi:10.1103/PhysRevD.91.124020.
  %[arXiv:1704.07065 [gr-qc]].

%\cite{Shaikh:2018kfv}
\bibitem{Shaikh:2018kfv}
  R.~Shaikh,
  Phys.\ Rev.\ D {\bf 98},  024044 (2018)
 % doi:10.1103/PhysRevD.98.024044.
  %[arXiv:1803.11422 [gr-qc]].

%\cite{Gyulchev:2018fmd}
\bibitem{Gyulchev:2018fmd}
  G.~Gyulchev, P.~Nedkova, V.~Tinchev and S.~Yazadjiev,
  %``On the shadow of rotating traversable wormholes,''
  Eur.\ Phys.\ J.\ C {\bf 78} ,  544 (2018)
  %doi:10.1140/epjc/s10052-018-6012-9.
  %[arXiv:1805.11591 [gr-qc]].

%\cite{Guerrero:2022qkh}
\bibitem{Guerrero:2022qkh}
M.~Guerrero, G.~J.~Olmo, D.~Rubiera-Garcia and D.~G\'omez S\'aez-Chill\'on,
%``Light ring images of double photon spheres in black hole and wormhole spacetimes,''
Phys. Rev. D \textbf{105}, no.8, 084057 (2022)
%doi:10.1103/PhysRevD.105.084057
%[arXiv:2202.03809 [gr-qc]].

%\cite{Harko:2008vy}
\bibitem{Harko:2008vy}
  T.~Harko, Z.~Kovacs and F.~S.~N.~Lobo,
  %``Electromagnetic signatures of thin accretion disks in wormhole geometries,''
  Phys.\ Rev.\ D {\bf 78}, 084005 (2008)
  %doi:10.1103/PhysRevD.78.084005.
  %[arXiv:0808.3306 [gr-qc]].

%\cite{Harko:2009xf}
\bibitem{Harko:2009xf}
  T.~Harko, Z.~Kovacs and F.~S.~N.~Lobo,
  %``Thin accretion disks in stationary axisymmetric wormhole spacetimes,''
  Phys.\ Rev.\ D {\bf 79}, 064001 (2009)
  %doi:10.1103/PhysRevD.79.064001.
  %[arXiv:0901.3926 [gr-qc]].

%\cite{Bambi:2013jda}
\bibitem{Bambi:2013jda}
  C.~Bambi,
  %``Broad Kα iron line from accretion disks around traversable wormholes,''
  Phys.\ Rev.\ D {\bf 87}, 084039 (2013)
  %doi:10.1103/PhysRevD.87.084039.
  %[arXiv:1303.0624 [gr-qc]].

%\cite{Zhou:2016koy}
\bibitem{Zhou:2016koy}
  M.~Zhou, A.~Cardenas-Avendano, C.~Bambi, B.~Kleihaus and J.~Kunz,
  %``Search for astrophysical rotating Ellis wormholes with X-ray reflection spectroscopy,''
  Phys.\ Rev.\ D {\bf 94}, 024036 (2016)
  %doi:10.1103/PhysRevD.94.024036.
  %[arXiv:1603.07448 [gr-qc]].

%\cite{Lamy:2018zvj}
\bibitem{Lamy:2018zvj}
  F.~Lamy, E.~Gourgoulhon, T.~Paumard and F.~H.~Vincent,
  %``Imaging a non-singular rotating black hole at the center of the Galaxy,''
  Class.\ Quant.\ Grav.\  {\bf 35},   115009 (2018)
 % doi:10.1088/1361-6382/aabd97.
  %[arXiv:1802.01635 [gr-qc]].

%\cite{Deligianni:2021ecz}
\bibitem{Deligianni:2021ecz}
E.~Deligianni, J.~Kunz, P.~Nedkova, S.~Yazadjiev and R.~Zheleva,
%``Quasiperiodic oscillations around rotating traversable wormholes,''
Phys. Rev. D \textbf{104},  024048 (2021)
%doi:10.1103/PhysRevD.104.024048.
%[arXiv:2103.13504 [gr-qc]].

%\cite{Deligianni:2021hwt}
\bibitem{Deligianni:2021hwt}
E.~Deligianni, B.~Kleihaus, J.~Kunz, P.~Nedkova and S.~Yazadjiev,
%``Quasiperiodic oscillations in rotating Ellis wormhole spacetimes,''
Phys. Rev. D \textbf{104},  064043 (2021)
%doi:10.1103/PhysRevD.104.064043.
%[arXiv:2107.01421 [gr-qc]].

%\cite{Damour:2007ap}
\bibitem{Damour:2007ap}
  T.~Damour and S.~N.~Solodukhin,
  %``Wormholes as black hole foils,''
  Phys.\ Rev.\ D {\bf 76}, 024016 (2007)
  %doi:10.1103/PhysRevD.76.024016
  %[arXiv:0704.2667 [gr-qc]].                     

%\cite{Azreg-Ainou:2014dwa}
\bibitem{Azreg-Ainou:2014dwa}
  M.~Azreg-Ainou,
  %``Confined-exotic-matter wormholes with no gluing effects—Imaging supermassive wormholes and black holes,''
  JCAP {\bf 1507}, 037 (2015)
  %doi:10.1088/1475-7516/2015/07/037
  %[arXiv:1412.8282 [gr-qc]].

%\cite{Dzhunushaliev:2016ylj}
\bibitem{Dzhunushaliev:2016ylj}
  V.~Dzhunushaliev, V.~Folomeev, B.~Kleihaus and J.~Kunz,
  %``Can mixed star-plus-wormhole systems mimic black holes?,''
  JCAP {\bf 1608}, no. 08, 030 (2016)
  %doi:10.1088/1475-7516/2016/08/030
  %[arXiv:1601.04124 [gr-qc]].

%\cite{Cardoso:2016rao}
\bibitem{Cardoso:2016rao}
  V.~Cardoso, E.~Franzin and P.~Pani,
  %``Is the gravitational-wave ringdown a probe of the event horizon?,''
  Phys.\ Rev.\ Lett.\  {\bf 116}, 171101 (2016);
  Erratum: [Phys.\ Rev.\ Lett.\  {\bf 117}, 089902 (2016)]
  %doi:10.1103/PhysRevLett.117.089902, 10.1103/PhysRevLett.116.171101
  %[arXiv:1602.07309 [gr-qc]].

%\cite{Konoplya:2016hmd}
\bibitem{Konoplya:2016hmd}
  R.~A.~Konoplya and A.~Zhidenko,
  %``Wormholes versus black holes: quasinormal ringing at early and late times,''
  JCAP {\bf 1612}, 043 (2016)
  %doi:10.1088/1475-7516/2016/12/043
  %[arXiv:1606.00517 [gr-qc]].

%\cite{Nandi:2016uzg}
\bibitem{Nandi:2016uzg}
  K.~K.~Nandi, R.~N.~Izmailov, A.~A.~Yanbekov and A.~A.~Shayakhmetov,
  %``Ring-down gravitational waves and lensing observables: How far can a wormhole mimic those of a black hole?,''
  Phys.\ Rev.\ D {\bf 95}, 104011 (2017)
  %doi:10.1103/PhysRevD.95.104011
  %[arXiv:1611.03479 [gr-qc]].

%\cite{Bueno:2017hyj}
\bibitem{Bueno:2017hyj}
  P.~Bueno, P.~A.~Cano, F.~Goelen, T.~Hertog and B.~Vercnocke,
  %``Echoes of Kerr-like wormholes,''
  Phys.\ Rev.\ D {\bf 97}, 024040 (2018)
  %doi:10.1103/PhysRevD.97.024040
  %[arXiv:1711.00391 [gr-qc]].

%\cite{Blazquez-Salcedo:2018ipc}
\bibitem{Blazquez-Salcedo:2018ipc}
  J.~L.~Bl\'{a}zquez-Salcedo, X.~Y.~Chew and J.~Kunz,
  %``Scalar and axial quasinormal modes of massive static phantom wormholes,''
  Phys.\ Rev.\ D {\bf 98} no.4,  044035 (2018)
  %doi:10.1103/PhysRevD.98.044035
  %[arXiv:1806.03282 [gr-qc]].

%\cite{Azad:2020ajs}
\bibitem{Azad:2020ajs}
B.~Azad, F.~Loran and A.~Mostafazadeh,
%``Transmission of low-energy scalar waves through a traversable wormhole,''
Eur. Phys. J. C \textbf{80} (2020) no.12, 1097
%doi:10.1140/epjc/s10052-020-08668-3
%[arXiv:2010.15023 [gr-qc]].
%3 citations counted in INSPIRE as of 12 Jan 2023

  %\cite{Azad:2022qqn}
\bibitem{Azad:2022qqn}
B.~Azad, J.~L.~Bl\'azquez-Salcedo, X.~Y.~Chew, J.~Kunz and D.~h.~Yeom,
%``Polar modes and isospectrality of Ellis-Bronnikov wormholes,''
[arXiv:2212.12601 [gr-qc]]

 %\cite{LIGOScientific:2016aoc}
\bibitem{LIGOScientific:2016aoc}
B.~P.~Abbott \textit{et al.} [LIGO Scientific and Virgo],
%``Observation of Gravitational Waves from a Binary Black Hole Merger,''
Phys. Rev. Lett. \textbf{116}, no.6, 061102 (2016)
%doi:10.1103/PhysRevLett.116.061102
%[arXiv:1602.03837 [gr-qc]].

%\cite{LIGOScientific:2017vwq}
\bibitem{LIGOScientific:2017vwq}
B.~P.~Abbott \textit{et al.} [LIGO Scientific and Virgo],
%``GW170817: Observation of Gravitational Waves from a Binary Neutron Star Inspiral,''
Phys. Rev. Lett. \textbf{119} , no.16, 161101 (2017)
%doi:10.1103/PhysRevLett.119.161101
%[arXiv:1710.05832 [gr-qc]].

%\cite{Kokkotas:1999bd}
\bibitem{Kokkotas:1999bd}
K.~D.~Kokkotas and B.~G.~Schmidt,
%``Quasinormal modes of stars and black holes,''
Living Rev. Rel. \textbf{2}, 2 (1999)
%doi:10.12942/lrr-1999-2
%[arXiv:gr-qc/9909058 [gr-qc]].

%\cite{Berti:2009kk}
\bibitem{Berti:2009kk}
E.~Berti, V.~Cardoso and A.~O.~Starinets,
%``Quasinormal modes of black holes and black branes,''
Class. Quant. Grav. \textbf{26}, 163001 (2009)
%doi:10.1088/0264-9381/26/16/163001
%[arXiv:0905.2975 [gr-qc]].

%\cite{Konoplya:2011qq}
\bibitem{Konoplya:2011qq}
R.~A.~Konoplya and A.~Zhidenko,
%``Quasinormal modes of black holes: From astrophysics to string theory,''
Rev. Mod. Phys. \textbf{83}, 793 (2011)
%doi:10.1103/RevModPhys.83.793
%[arXiv:1102.4014 [gr-qc]].


%\cite{Ellis:1973yv}
\bibitem{Ellis:1973yv} 
  H.~G.~Ellis,
  %``Ether flow through a drainhole - a particle model in general relativity,''
  J.\ Math.\ Phys.\  {\bf 14}, 104 (1973)
 % doi:10.1063/1.1666161.

%\cite{Ellis:1979bh}
\bibitem{Ellis:1979bh} 
  H.~G.~Ellis,
  %``The Evolving, Flowless Drain Hole: A Nongravitating Particle Model In General Relativity Theory,''
  Gen.\ Rel.\ Grav.\  {\bf 10}, 105 (1979)
 % doi:10.1007/BF00756794.

%\cite{Bronnikov:1973fh}
\bibitem{Bronnikov:1973fh} 
  K.~A.~Bronnikov,
  %``Scalar-tensor theory and scalar charge,''
  Acta Phys.\ Polon.\ B {\bf 4}, 251 (1973)

%\cite{Morris:1988tu}
\bibitem{Morris:1988tu}
M.~S.~Morris, K.~S.~Thorne and U.~Yurtsever,
%``Wormholes, Time Machines, and the Weak Energy Condition,''
Phys. Rev. Lett. \textbf{61}, 1446-1449 (1988)
%doi:10.1103/PhysRevLett.61.1446

%\cite{Morris:1988cz}
\bibitem{Morris:1988cz}
M.~S.~Morris and K.~S.~Thorne,
%``Wormholes in space-time and their use for interstellar travel: A tool for teaching general relativity,''
Am. J. Phys. \textbf{56}, 395 (1988)
%doi:10.1119/1.15620

%\cite{Visser:1995cc}
\bibitem{Visser:1995cc} 
  M.~Visser,
%  ``Lorentzian wormholes: From Einstein to Hawking,''
  Woodbury, USA: AIP (1995)

%\cite{Hochberg:1996ee}
\bibitem{Hochberg:1996ee}
D.~Hochberg, A.~Popov and S.~V.~Sushkov,
%``Selfconsistent wormhole solutions of semiclassical gravity,''
Phys. Rev. Lett. \textbf{78}, 2050 (1997)
%doi:10.1103/PhysRevLett.78.2050
%[arXiv:gr-qc/9701064 [gr-qc]].

%\cite{Hochberg:1998ii}
\bibitem{Hochberg:1998ii}
D.~Hochberg and M.~Visser,
%``The Null energy condition in dynamic wormholes,''
Phys. Rev. Lett. \textbf{81}, 746 (1998)
%doi:10.1103/PhysRevLett.81.746
%[arXiv:gr-qc/9802048 [gr-qc]].

%\cite{Visser:2003yf}
\bibitem{Visser:2003yf}
M.~Visser, S.~Kar and N.~Dadhich,
%``Traversable wormholes with arbitrarily small energy condition violations,''
Phys. Rev. Lett. \textbf{90}, 201102 (2003)
%doi:10.1103/PhysRevLett.90.201102
%[arXiv:gr-qc/0301003 [gr-qc]].

%\cite{Bronnikov:2005gm}
\bibitem{Bronnikov:2005gm}
K.~A.~Bronnikov and J.~C.~Fabris,
%``Regular phantom black holes,''
Phys. Rev. Lett. \textbf{96}, 251101 (2006)
%doi:10.1103/PhysRevLett.96.251101
%[arXiv:gr-qc/0511109 [gr-qc]].

%\cite{Alcubierre:2017pqm}
\bibitem{Alcubierre:2017pqm}
M.~Alcubierre and F.~S.~N.~Lobo,
%``Wormholes, Warp Drives and Energy Conditions,''
Fundam. Theor. Phys. \textbf{189} (2017) 
%doi:10.1007/978-3-319-55182-1.
%[arXiv:2103.05610 [gr-qc]].

%\cite{Kanti:2011jz}
\bibitem{Kanti:2011jz}
P.~Kanti, B.~Kleihaus and J.~Kunz,
%``Wormholes in Dilatonic Einstein-Gauss-Bonnet Theory,''
Phys. Rev. Lett. \textbf{107}, 271101 (2011)
%doi:10.1103/PhysRevLett.107.271101
%[arXiv:1108.3003 [gr-qc]].

%\cite{Blazquez-Salcedo:2020czn}
\bibitem{Blazquez-Salcedo:2020czn}
J.~L.~Bl\'azquez-Salcedo, C.~Knoll and E.~Radu,
%``Traversable wormholes in Einstein-Dirac-Maxwell theory,''
Phys. Rev. Lett. \textbf{126}, 101102 (2021)
%doi:10.1103/PhysRevLett.126.101102
%[arXiv:2010.07317 [gr-qc]].

%\cite{Konoplya:2021hsm}
\bibitem{Konoplya:2021hsm}
R.~A.~Konoplya and A.~Zhidenko,
%``Traversable Wormholes in General Relativity,''
Phys. Rev. Lett. \textbf{128}, 091104 (2022)
%doi:10.1103/PhysRevLett.128.091104
%[arXiv:2106.05034 [gr-qc]].

%\cite{Shinkai:2002gv}
\bibitem{Shinkai:2002gv}
  H.~a.~Shinkai and S.~A.~Hayward,
  %``Fate of the first traversible wormhole: Black hole collapse or inflationary expansion,''
  Phys.\ Rev.\ D {\bf 66}, 044005 (2002)
%  doi:10.1103/PhysRevD.66.044005
%  [gr-qc/0205041].

%\cite{Gonzalez:2008wd}
\bibitem{Gonzalez:2008wd}
  J.~A.~Gonzalez, F.~S.~Guzman and O.~Sarbach,
  %``Instability of wormholes supported by a ghost scalar field. I. Linear stability analysis,''
  Class.\ Quant.\ Grav.\  {\bf 26}, 015010 (2009)
 % doi:10.1088/0264-9381/26/1/015010
%  [arXiv:0806.0608 [gr-qc]].

%\cite{Gonzalez:2008xk}
\bibitem{Gonzalez:2008xk}
  J.~A.~Gonzalez, F.~S.~Guzman and O.~Sarbach,
  %``Instability of wormholes supported by a ghost scalar field. II. Nonlinear evolution,''
  Class.\ Quant.\ Grav.\  {\bf 26},  015011 (2009)
 % doi:10.1088/0264-9381/26/1/015011
%  [arXiv:0806.1370 [gr-qc]].

  %\cite{Cremona:2018wkj}
\bibitem{Cremona:2018wkj}
F.~Cremona, F.~Pirotta and L.~Pizzocchero,
%``On the linear instability of the Ellis\textendash{}Bronnikov\textendash{}Morris-Thorne wormhole,''
Gen. Rel. Grav. \textbf{51}, 19 (2019)
%doi:10.1007/s10714-019-2501-x
%[arXiv:1805.02602 [gr-qc]].

%\cite{Kashargin:2007mm}
\bibitem{Kashargin:2007mm}
P.~E.~Kashargin and S.~V.~Sushkov,
%``Slowly rotating wormholes: The First order approximation,''
Grav. Cosmol. \textbf{14}, 80 (2008)
%doi:10.1134/S0202289308010106
%[arXiv:0710.5656 [gr-qc]].

%\cite{Kashargin:2008pk}
\bibitem{Kashargin:2008pk}
P.~E.~Kashargin and S.~V.~Sushkov,
%``Slowly rotating scalar field wormholes: The Second order approximation,''
Phys. Rev. D \textbf{78}, 064071 (2008)
%doi:10.1103/PhysRevD.78.064071
%[arXiv:0809.1923 [gr-qc]].

%\cite{Kleihaus:2014dla}
\bibitem{Kleihaus:2014dla}
B.~Kleihaus and J.~Kunz,
%``Rotating Ellis Wormholes in Four Dimensions,''
Phys. Rev. D \textbf{90}, 121503 (2014)
%doi:10.1103/PhysRevD.90.121503
%[arXiv:1409.1503 [gr-qc]].

%\cite{Chew:2016epf}
\bibitem{Chew:2016epf}
X.~Y.~Chew, B.~Kleihaus and J.~Kunz,
%``Geometry of Spinning Ellis Wormholes,''
Phys. Rev. D \textbf{94}, 104031 (2016)
%doi:10.1103/PhysRevD.94.104031
%[arXiv:1608.05253 [gr-qc]].

%\cite{Dzhunushaliev:2013jja}
\bibitem{Dzhunushaliev:2013jja}
V.~Dzhunushaliev, V.~Folomeev, B.~Kleihaus, J.~Kunz and E.~Radu,
%``Rotating Wormholes in Five Dimensions,''
Phys. Rev. D \textbf{88}, 124028 (2013)
%doi:10.1103/PhysRevD.88.124028
%[arXiv:1309.2448 [gr-qc]].

%\cite{Dias:2015wqa}
\bibitem{Dias:2015wqa}
O.~J.~C.~Dias, M.~Godazgar and J.~E.~Santos,
%``Linear Mode Stability of the Kerr-Newman Black Hole and Its Quasinormal Modes,''
Phys. Rev. Lett. \textbf{114}, 151101 (2015)
%doi:10.1103/PhysRevLett.114.151101
%[arXiv:1501.04625 [gr-qc]].

%\cite{Blazquez-Salcedo:2022eik}
\bibitem{Blazquez-Salcedo:2022eik}
J.~L.~Bl\'azquez-Salcedo and F.~S.~Khoo,
%``Quasinormal modes of slowly rotating Kerr-Newman black holes using the double series method,''
[arXiv:2212.00054 [gr-qc]]

%\cite{Azad:2023}
\bibitem{Azad:2023}
B.~Azad, J.~L.~Bl\'azquez-Salcedo, F.~S.~Khoo and J.~Kunz,
in preparation

%\cite{Matos:2005uh}
\bibitem{Matos:2005uh}
T.~Matos and D.~Nunez,
%``Rotating scalar field wormhole,''
Class. Quant. Grav. \textbf{23}, 4485 (2006)
%doi:10.1088/0264-9381/23/13/012
%[arXiv:gr-qc/0508117 [gr-qc]].

\end{thebibliography}


%%%%%%%%%%%%%%%%%%%%%%%%%%%%%%%%%%%%%%%%%%%%%%%%%%%%%%%%%%%%%%%%%%%%%%%%%


\end{document}
 


